\section{Prepare and observe the cumulative program}
Start from \texttt{drafts/crc-construction/main.asm}. Replace one named marked region
at a time, retaining the surrounding labels. The first bit/table stubs contain three
instructions because they establish a defined incoming-state return before branching.
Other marked stubs contain a branch and its delay-slot \texttt{nop}. The initializer's
zero return and the file routine's development status three are provisional markers.
They do not establish that all requested work is complete.

The selector word \texttt{exercise} chooses bit update zero, initializer one, table
update two or file operation three. Keep selector, count, seed, capacity and method at
their supplied defaults in the cumulative source. The preparation helper changes them
in a separate observation file. For table-update observations it also supplies an
independently computed table, so those checks do not rely on the student's initializer.

Use \texttt{prepare.py} with a recipe ID, \texttt{--source} pointing to the current
cumulative source, and a new \texttt{--output} path. The source is never overwritten.
For a real file observation, optional \texttt{--fixture-output} writes the recipe's
binary bytes to another new path and embeds its absolute ASCII path in the observation
source. The path must fit the supplied 128-byte arena. These files are ordinary host
files; they do not force short reads or injected errors. The CLI guided checks use
explicit in-memory fixture controls for those observations.

Most kernel observations stop at \texttt{kernel\_done} before the caller stores the
raw result. File observations stop at \texttt{file\_done} before its two stores. Direct
checks supply \texttt{function\_return} instead, and also use nonzero source, table
and scratch offsets. The file internal observation stops on the second visit to
\texttt{file\_read}; working saved registers are inspected there, before restoration.
Continue the same session to check the completed return and descriptor cleanup.

Unfinished file regions branch to a supplied cleanup path. The frame initializes a
negative descriptor marker, so cleanup closes only a successfully acquired descriptor.
It returns the retained raw state and development status three, then restores the
frame. The completed success and error paths bypass this temporary path. Final
interface checks never accept development status three as a completed file result.
